\section*{Chapitre \ref{ch:b2:microbial-metabolism} --- Métabolismes microbiens et biofilms}

\begin{solution}{exo:b2:microbial-metabolism:1}
Cyanobactérie~: photo-litho-autotrophe (lumière, eau, $\mathrm{CO_2}$).
\emph{Nitrosomonas}~: chimio-litho-autotrophe (oxydation de l'ammonium,
$\mathrm{CO_2}$). \emph{E.~coli} sur glucose~: chimio-organo-hétérotrophe.
Bactérie pourpre non sulfureuse sur succinate à la lumière~:
photo-organo-hétérotrophe. \emph{Thiobacillus} sur sulfure à l'obscurité~:
chimio-litho-autotrophe.
\end{solution}

\begin{solution}{exo:b2:microbial-metabolism:2}
Le dioxygène (millimètres de surface), le nitrate (juste en dessous, là
où le dioxygène a disparu), les oxydes de manganèse et de fer (la
transition du brun au gris), le sulfate (la couche noire sulfurée,
épaisse dans les sédiments marins), le
$\mathrm{CO_2}$ pour la \omterm{def:b2:microbial-metabolism:anaerobic}{méthanogenèse} (au plus profond, ou juste sous la
surface dans une eau douce pauvre en sulfate).
\end{solution}

\begin{solution}{exo:b2:microbial-metabolism:3}
$5\,\mathrm{C_6H_{12}O_6} + 24\,\mathrm{NO_3^-} + 24\,\mathrm{H^+}
\to 30\,\mathrm{CO_2} + 12\,\mathrm{N_2} + 42\,\mathrm{H_2O}$~: 4.8
ions nitrate par glucose (chaque nitrate accepte 5 électrons, chaque
glucose en cède 24).
\end{solution}

\begin{solution}{exo:b2:microbial-metabolism:4}
Une communauté microbienne fixée à une surface et enrobée dans une
matrice qu'elle produit elle-même, faite de polysaccharides, de
protéines et d'ADN. Étapes~: adhésion, microcolonie et sécrétion de
matrice, maturation avec tours et canaux, dispersion. La plaque
dentaire et les caries~; les infections des cathéters et des implants
(et des poumons dans la mucoviscidose)~; l'encrassement et la corrosion
des canalisations et des coques --- ou, utilement, les lits bactériens
du traitement des eaux.
\end{solution}

\begin{solution}{exo:b2:microbial-metabolism:5}
Huit électrons descendent de \qty{-0.42}{V} à \qty{-0.24}{V}~:
$\Delta G^{\circ\prime} = -8\times 96.5\times 0.18 =
\qty{-139}{kJ/mol}$ de méthane (valeur tabulée \qty{-131}{kJ/mol}),
soit \qty{35}{kJ} par $\mathrm{H_2}$~: au plus 0.7 ATP par dihydrogène,
en pratique environ un demi. L'énergie par molécule de substrat est si
faible que les méthanogènes croissent lentement (\omterm{def:b2:unicellular-diversity:division}{temps de génération}
de quelques heures à quelques jours) et convertissent presque tout leur
substrat en gaz plutôt qu'en cellules.
\end{solution}

\begin{solution}{exo:b2:microbial-metabolism:6}
$\mu = 1.2\,S/(5 + S)$~: \qty{0.20}{h^{-1}} ($T = \ln 2/\mu =
\qty{3.5}{h}$), \qty{0.60}{h^{-1}} (\qty{1.2}{h}), \qty{1.09}{h^{-1}}
(\qty{0.64}{h}). Quatre-vingt-dix pour cent~: $S = 9K_s = \qty{45}{mg/L}$.
\end{solution}

\begin{solution}{exo:b2:microbial-metabolism:7}
$S^* = 0.05\times 0.4/(0.8 - 0.4) = \qty{0.05}{g/L}$~; $X^* =
0.4\times(5 - 0.05) = \qty{1.98}{g/L}$~; production $DX^* = 0.4\times
1.98 = \qty{0.79}{g/L/h}$~; lessivage $D_c = 0.8\times 5/5.05 =
\qty{0.79}{h^{-1}}$.
\end{solution}

\begin{solution}{exo:b2:microbial-metabolism:8}
Énergie captée~: $0.25\times 275 = \qty{69}{kJ}$, soit 1.4 ATP par
ammonium~; $12/1.4 = 8.7$, environ neuf ions ammonium par carbone ---
et plusieurs fois plus une fois compté le coût de la production du NADH
par remontée des électrons depuis le nitrite. Un hétérotrophe reçoit son
carbone déjà réduit et organisé, et trente ATP par glucose~; un
nitrifiant doit acheter à la fois son énergie et son pouvoir réducteur à
un donneur médiocre et construire chaque carbone à partir du
$\mathrm{CO_2}$, si bien qu'il convertit en cellules une infime fraction
du substrat qu'il traite.
\end{solution}

\begin{solution}{exo:b2:microbial-metabolism:9}
$L_p = \sqrt{2\times 1.2\times 10^{-9}\times 0.2/0.5} =
\sqrt{9.6\times 10^{-10}} = \qty{31}{\micro m}$. Doubler $c_0$
multiplie $L_p$ par $\sqrt 2$~: \qty{44}{\micro m}~; quadrupler $q$ la
divise par deux~: \qty{15}{\micro m}.
\end{solution}

\begin{solution}{exo:b2:microbial-metabolism:10}
De haut en bas~: l'eau avec algues et cyanobactéries (photosynthèse
oxygénique)~; une bande couleur rouille de ferro-oxydants là où
$\mathrm{Fe^{2+}}$ rencontre le dioxygène~; un voile blanc de
sulfo-oxydants incolores (\omterm{def:b2:microbial-metabolism:trophic}{chimiolithotrophes}) là où le sulfure rencontre
le dioxygène~; une bande pourpre de bactéries pourpres sulfureuses
(\omterm{def:b2:microbial-metabolism:anoxygenic}{photosynthèse anoxygénique} sur sulfure)~; une bande verte de bactéries
vertes sulfureuses (même métabolisme, tolérant plus de sulfure et moins
de lumière)~; la vase noire des fermenteurs et des sulfato-réducteurs.
(a) Les bactéries pourpres sulfureuses ont besoin à la fois de la
lumière venue d'en haut et du sulfure venu d'en bas~; elles se placent
donc à l'interface où les deux gradients se croisent. (b) Les
sulfato-réducteurs produisent $\mathrm{H_2S}$, qui précipite le fer sous
forme de FeS noir. (c) Fermée pour la matière~: soufre, carbone et azote
circulent entre formes oxydées et réduites sans sortir~; ouverte pour
l'énergie~: la lumière entre, la chaleur sort, et sans lumière tous les
cycles s'arrêtent.
\end{solution}

\begin{solution}{exo:b2:microbial-metabolism:11}
Production $pn$, perte $kA$~: $A^* = pn/k$. $A_c = 10^{-8}\times
6.02\times 10^{23} = 6.0\times 10^{15}$ molécules par litre~; $n_c =
kA_c/p = 5\times 10^{-4}\times 6.0\times 10^{15}/5 = 6\times 10^{11}$
par litre, soit $6\times 10^{8}$ par millilitre. La culture dense
($10^{9}$) dépasse le seuil, l'eau de mer ($10^{5}$) est quatre ordres
de grandeur en dessous. Dans un \omterm{def:b2:microbial-metabolism:biofilm}{biofilm}, les cellules sont à $10^{11}$
par millilitre de matrice, et la matrice ralentit la fuite du signal
(ce qui revient à abaisser $k$)~: une microcolonie de quelques milliers
de cellules dans un nanolitre atteint le quorum, alors que les mêmes
cellules dispersées dans un litre ne l'atteindraient jamais.
\end{solution}

\begin{solution}{exo:b2:microbial-metabolism:12}
La fixation de l'azote (seuls les \omterm{def:b2:unicellular-diversity:prokaryote}{procaryotes} possèdent la
nitrogénase)~: si elle cessait, l'azote combiné de la biosphère
s'écoulerait par \omterm{def:b2:microbial-metabolism:anaerobic}{dénitrification} et la production s'effondrerait en
quelques décennies. La nitrification et la \omterm{def:b2:microbial-metabolism:anaerobic}{dénitrification} (\omterm{def:b2:unicellular-diversity:prokaryote}{procaryotes}
seulement)~: si elles cessaient, l'ammonium s'accumulerait, le nitrate
disparaîtrait et l'azote ne retournerait jamais à l'air. La
\omterm{def:b2:microbial-metabolism:anaerobic}{méthanogenèse} et la \omterm{def:b2:microbial-metabolism:anaerobic}{respiration anaérobie} (\omterm{def:b2:unicellular-diversity:prokaryote}{procaryotes} seulement)~: si
elles cessaient, la matière organique des sédiments anoxiques, des
panses et des marais s'accumulerait sans être minéralisée et le carbone
sortirait du cycle~; de même, la photosynthèse oxygénique fut une
invention bactérienne, et sans elle il n'y aurait pas de dioxygène à
respirer. Les eucaryotes sont un ajout tardif et chimiquement limité à
une planète \omterm{def:b2:unicellular-diversity:prokaryote}{procaryote}.
\end{solution}

\begin{solution}{pb:b2:microbial-metabolism:1}
\textbf{1.} $-24\times 96.5\times (0.82 + 0.43) = \qty{-2900}{kJ/mol}$.
\textbf{2.} Nitrate~: $-24\times 96.5\times 1.17 = \qty{-2700}{kJ/mol}$~;
sulfate~: $-24\times 96.5\times 0.21 = \qty{-490}{kJ/mol}$.
\textbf{3.} Le dioxygène rapporte davantage, si bien que les aérobies
supplantent les dénitrifiants pour la matière organique tant qu'il
reste du dioxygène, et le dioxygène réprime les nitrate réductases~; la
\omterm{def:b2:microbial-metabolism:anaerobic}{sulfato-réduction} exige l'anoxie \emph{et} l'absence de nitrate~: son
odeur signifie donc que le bassin est sous-aéré et le nitrate épuisé
--- et $\mathrm{H_2S}$ est toxique et corrode le béton.
\textbf{4.} $0.4\times 2900/50 = 23$ ATP~; $0.4\times 2700/50 = 22$~;
$0.4\times 490/50 = 4$.
\textbf{5.} Environ six fois plus ($23/4$).
\textbf{6.} Deux ATP par glucose correspondent à un rendement d'environ
\num{0.1}~: \qty{90}{\%} de l'énergie du substrat restent dans les
produits (acides, alcools, $\mathrm{H_2}$), que les méthanogènes
convertissent en méthane avec un gain tout aussi faible~; peu d'énergie,
peu de biomasse, et le carbone part sous forme de gaz.
\textbf{7.} $10^4\times 0.3 = \qty{3000}{kg}$ de DCO par jour.
\textbf{8.} Biomasse $0.4\times 3000 = \qty{1200}{kg}$ de DCO~; les
\qty{1800}{kg} de DCO restants sont oxydés et exigent \qty{1800}{kg}
de $\mathrm{O_2}$.
\textbf{9.} $1800/2 = \qty{900}{kWh}$ par jour.
\textbf{10.} $10^4\times 0.04 = \qty{400}{kg}$ d'azote~;
$4.57\times 400 = \qty{1830}{kg}$ de $\mathrm{O_2}$.
\textbf{11.} En régime stationnaire, $\mu = 1/\theta \le \mu_{\max}$~:
lessivage en dessous de $\theta = 1/\mu_{\max} = \qty{2}{d}$.
\textbf{12.} $D = \qty{0.1}{d^{-1}}$~: $S^* = 1\times 0.1/(0.5 - 0.1)
= \qty{0.25}{g/m^3}$.
\textbf{13.} $S^* = 0.1/(0.25 - 0.1) = \qty{0.67}{g/m^3}$, et la limite
de lessivage passe à \qty{4}{d}~; l'exploitant doit allonger l'âge des
boues (en extraire moins)~: à $\theta = \qty{20}{d}$, $S^* =
0.05/0.2 = \qty{0.25}{g/m^3}$ de nouveau.
\textbf{14.} $2.86\times 400 = \qty{1144}{kg}$ de DCO par jour.
\textbf{15.} Ammonium de l'effluent $0.25\times 10^4 = \qty{2.5}{kg/d}$,
soit \qty{0.6}{\%} de la charge~; le reste, \qty{397}{kg} d'azote, part
sous forme de $\mathrm{N_2}$~: \qty{99}{\%}.
\textbf{16.} Économie~: \qty{1144}{kg} de $\mathrm{O_2}$. Demande
totale~: $(1800 - 1144) + 1830 = \qty{2490}{kg}$ de $\mathrm{O_2}$ par
jour.
\textbf{17.} $1200\times 0.6\times 0.35 = \qty{252}{m^3}$ de méthane
par jour.
\textbf{18.} $252\times 36 = \qty{9070}{MJ} = \qty{2520}{kWh}$~;
électricité $0.35\times 2520 = \qty{880}{kWh}$ par jour.
\textbf{19.} Aération $2490/2 = \qty{1245}{kWh}$ par jour~: le méthane
en couvre $880/1245 = \qty{71}{\%}$.
\textbf{20.} $D_e\,c'' = q$, donc $c = c_0 + ax + qx^2/2D_e$~; avec
$c(L_p) = 0$ et $c'(L_p) = 0$~: $a = -qL_p/D_e$ et $c =
(q/2D_e)(L_p - x)^2$, avec $L_p^2 = 2D_ec_0/q$.
\textbf{21.} $L_p = \sqrt{2\times 1.5\times 10^{-9}\times 0.2/0.3} =
\sqrt{2\times 10^{-9}} = \qty{45}{\micro m}$.
\textbf{22.} $45/300 = \qty{15}{\%}$ du film est aérobie. En dessous,
des dénitrifiants réduisent en $\mathrm{N_2}$ le nitrate qui diffuse
depuis la couche nitrifiante de surface, en utilisant la matière
organique qui diffuse jusqu'à eux.
\textbf{23.} Dans un même film, le gradient de dioxygène crée une
couche aérobie (nitrification) au-dessus d'une couche anoxique
(\omterm{def:b2:microbial-metabolism:anaerobic}{dénitrification}), à quelques dizaines de micromètres l'une de l'autre~;
le bassin aéré est oxique partout, si bien que la \omterm{def:b2:microbial-metabolism:anaerobic}{dénitrification} exige
un bassin anoxique séparé.
\textbf{24.} Le chlore réagit avec les polymères de la matrice et y est
consommé avant d'atteindre les cellules profondes~; les cellules
profondes, affamées, à croissance lente ou dormantes, sont bien moins
sensibles~; seule la couche superficielle meurt et le film se reforme
à partir du dessous.
\textbf{25.} Demande en dioxygène \qty{2490}{kg} par jour~; méthane
\qty{252}{m^3} par jour~; le méthane fournit \qty{71}{\%} de
l'électricité d'aération (\qty{880}{kWh} sur \qty{1245}{kWh}).
\end{solution}
